% Figure: an isolated joint free body, showing the two member forces
% meeting at the joint and the external load, with the equilibrium
% triangle drawn alongside.
\begin{figure}[htbp]
\centering
\begin{tikzpicture}[>=latex, scale=1.0]
  % === left panel: the joint with members ===
  \begin{scope}
    \coordinate (J) at (0,0);
    % horizontal member to the right (chord)
    \draw[engblue, very thick] (J) -- (2.2,0);
    % diagonal member down-right at 45 deg
    \draw[engblue, very thick] (J) -- (-1.6,-1.6);
    \fill[engblack] (J) circle (2pt);
    \node[anchor=south] at (J) {joint $J$};
    % member force arrows (assumed tension pulls away from joint)
    \draw[->, thick, engred] (J) -- (1.4,0);
    \node[engred, anchor=south] at (1.0,0.05) {$F_1$};
    \draw[->, thick, engred] (J) -- (-1.0,-1.0);
    \node[engred, anchor=east] at (-0.9,-0.85) {$F_2$};
    % external load straight down
    \draw[->, very thick, engamber] (J) -- (0,-1.5);
    \node[engamber, anchor=west] at (0.05,-1.1) {$P$};
    % angle mark
    \draw[enggray, thin] (0.5,0) arc (0:225:0.5);
    \node[enggray, font=\scriptsize] at (-0.55,-0.25) {$\theta$};
  \end{scope}
  % === right panel: closed force polygon ===
  \begin{scope}[xshift=5.2cm, yshift=-0.8cm]
    % P down, F1 along +x, F2 along direction (-1,-1)/sqrt2; closed triangle
    \coordinate (O) at (0,0);
    \coordinate (A) at (0,-1.5);            % P (down)
    \coordinate (B) at (1.5,-1.5);          % then F1 horizontal closing? draw closed
    % closed polygon: O -> down (P) -> right (F1) -> back to O along F2
    \draw[->, very thick, engamber] (O) -- (A);
    \node[engamber, anchor=east] at ($(O)!0.5!(A)$) {$P$};
    \draw[->, thick, engred] (A) -- (B);
    \node[engred, anchor=north] at ($(A)!0.5!(B)$) {$F_1$};
    \draw[->, thick, engred] (B) -- (O);
    \node[engred, anchor=south west] at ($(B)!0.5!(O)$) {$F_2$};
    \node[enggray, font=\scriptsize, anchor=west] at (1.7,-0.7)
      {closed polygon $\Rightarrow \sum\vec{F}=\vec 0$};
  \end{scope}
\end{tikzpicture}
\caption[Method of joints: isolated joint and its force
polygon]{Method of joints. Left: a single joint $J$ cut free as a point
mass, carrying the axial forces $F_1$ and $F_2$ of the two members
meeting there and the external load $P$. By convention the unknown
member forces are drawn as tension, pulling away from the joint; a
negative answer then reports compression with no change of bookkeeping.
The joint supplies two scalar equations, $\sum F_x=0$ and $\sum F_y=0$,
exactly enough to solve a joint where only two unknown members meet.
Right: the same statement read graphically. Three forces in equilibrium
close into a triangle, the force polygon, whose closure is the
equilibrium condition drawn rather than written. The graphical and
algebraic methods give the same numbers; the polygon is the faster
sanity check on a sign.}
\label{fig:vol03ch02:method-of-joints}
\end{figure}
